\section{Saved predictive results}
\label{a2:results}

Both follow-ups used the same pinned 913-event GeoNet selection and three
chronological prediction windows containing 203, 155 and 254 events.
The three settings reuse those 612 distinct targets; they do not create
1,836 independent observations or nine independent trials. Parameters are
fitted using each training prefix and then held fixed. Genuinely observed
past events update the forecast history, with future events excluded.
The reported quantities are event-weighted log-score differences over the
complete forecast exposure, including quiet periods.

\subsection{The class-history extension}

Saved source records report nine observed cases, 30 fitting procedures
and all 240 required starts for the completed history study. They record
a passed independent saved-output audit, with no unknown procedure or
forecast-ambiguity flag. The supplied audit/start evidence excerpt retains
these completion records. The present editorial reconciliation checks
saved records and aggregates rather than repeating that audit. These are
completed local searches, not global-optimum certificates
\citep{MagnitudePredictionPilot2026}.

All nine constant-background fits prescribe \(d=0\), and the three
primary \(E1_{h=0}\) ablations prescribe \(h=0\). The middle-window
\(I0\) fit also reaches the estimated-coordinate boundary \(\eta=0\).
In post-1980 fold 0, \(I0\) reaches \(d=2\) and \(E1\) reaches
\(\delta=-1\). These boundary-dependent local fits do not establish
interior identification of the corresponding coordinates.

% Exact cells: history/primary/E1-E0/{timing,mark,joint}, etc.
\begin{table}[tbp]
\centering\small
\caption{History-study forecast gains in nats per event. Positive values
favor the first model. Post-1980 changes the training start while
retaining earlier parents. The \(h=0\) ablation was required in primary
only. All settings reuse the same 612 targets.}
\label{a2:history-results}
\begin{tabular}{llrrrr}
\toprule
Setting & Contrast & Timing & Mark & Joint & Winning windows\\
\midrule
Primary & \(E1-E0\) & -0.003092 & -0.005590 & -0.008682 & 1/3\\
Primary & \(E1-I0\) & +0.026624 & -0.005590 & +0.021034 & 2/3\\
Primary & \(E1-E1_{h=0}\) & +0.000242 & -0.005590 & -0.005348 & 1/3\\
Constant background & \(E1-E0\) & -0.005894 & -0.005602 & -0.011496 & 1/3\\
Constant background & \(E1-I0\) & +0.015408 & -0.005602 & +0.009806 & 3/3\\
Post-1980 & \(E1-E0\) & -0.006146 & -0.002550 & -0.008696 & 1/3\\
Post-1980 & \(E1-I0\) & +0.025900 & -0.002550 & +0.023350 & 2/3\\
\bottomrule
\end{tabular}
\end{table}

The full \(E1\) candidate beat \(I0\) but lost to \(E0\) in every aggregate
setting. Its primary association-specific gain relative to
\(E1_{h=0}\) was also negative: a small timing difference of
\(0.000242\) nats/event was outweighed by a mark-score loss of
\(0.005590\). Joint estimation can alter fitted effective scale and
fertility, so the nonzero timing difference does not contradict the
conditional ground-intensity identity.

The direct baseline contrast \(E0-I0\), obtained as
\((E1-I0)-(E1-E0)\), was positive in all three settings:
\(0.029716\), \(0.021302\) and \(0.032046\) nats/event, respectively.
Beating the inhibitory candidate therefore does not establish an added
benefit from class history. The frozen criteria required at least
\(0.01\) nats/event over both baselines and at least two shared winning
windows in primary, with the declared sensitivity requirements. Those
conditions failed, yielding \texttt{stop\_no\_robust\_joint\_gain}.

\subsection{Continuous-productivity sensitivity}

The productivity study's saved source records report nine cases, 36 main
procedures, 42 primary profile procedures and all 456 mandatory starts.
They record an independent auditor that reconstructed geometry from the
raw local catalogue and checked 1,116 train/fine candidate evaluations,
288 near-candidate forecasts and 42 descriptive profile forecasts. The
supplied audit/start evidence excerpt preserves the relevant saved counts
and completion records. The editorial reconciliation checks these saved
records and aggregates; it does not repeat the raw-data audit
\citep{ProductivitySlopePilot2026}.

% Exact cells: productivity/<setting>/<contrast>/{timing,mark,joint}.
\begin{table}[tbp]
\centering\small
\caption{Free-productivity forecast gains, in nats per event. The mark-score
differences are exactly zero, so the joint differences concern timing
(up to floating-point summation). Positive values favor the first model.}
\label{a2:productivity-results}
\begin{tabular}{lrrrr}
\toprule
Setting & \shortstack{E0\\free - fixed} &
\shortstack{I0\\free - fixed} &
\shortstack{E0 free\\- I0 free} &
\shortstack{E0 positive\\windows}\\
\midrule
Primary & -0.008782 & -0.009105 & +0.030039 & 1/3\\
Constant background & -0.010447 & -0.008689 & +0.019544 & 1/3\\
Post-1980 & -0.060929 & -0.054066 & +0.025183 & 0/3\\
\bottomrule
\end{tabular}
\end{table}

The primary \(E0\) estimates were
\(r=0.471908,\ 0.605853,\ 0.610446\). Across both families and all settings,
16 of 18 estimates were below \(0.7\); the two middle-window post-1980
estimates were above it. All 18 estimated \(r\) values were interior.
Other estimated coordinates reached constraints: \(I0_{\rm fixed}\)
and \(I0_{\rm free}\) both reached \(\eta=0\) in the middle
constant-background window and \(d=2\) in post-1980 fold 0. Thus the
productivity study has four such estimated-boundary fits, distinct from
the history-study cases above. The 12 prescribed constant-background \(d=0\) coordinates and 18 fixed
\(r=0.7\) coordinates are restrictions rather than estimated-boundary
findings.
These results do not establish interior identification of every coordinate.
Both nested free-\(r\) families improved training fit but lost forecast
score in every aggregate setting. Training improvement is partly guaranteed by
admitting the exact fixed-\(r\) floor, and is not evidence of a correct
physical slope.

The primary fixed-\(r\) profiles used the frozen grid
\(0,\ 0.2,\ 0.35,\ 0.5,\ 0.7,\ 0.85,\ 0.95\). For both families in the
first window, sampled nodes \(0.35\) and \(0.5\) lay within two nats of
the selected optimum. In the other four family/window profiles, no
sampled node lay within that band of the off-grid optimum. These are
sparse descriptive support sets, not confidence intervals or evidence
of precise identification. The maximum saved forecast spread among
near-training-equivalent candidates was approximately
\(9.04\times10^{-5}\) nats/event, below the frozen \(0.01\) diagnostic
threshold. It describes the saved candidate pool, not global uncertainty.

The stage-one continuation gate required primary \(E0\) gain of at least
\(0.01\) nats/event with at least two positive windows, plus positive
aggregate gains and at least two positive windows in both sensitivities,
complete audit and no forecast ambiguity. It failed while the resource
reserve remained available. The checkpoint was
\texttt{stop\_no\_robust\_productivity\_gain}; the free-\(r\) class-history
models \(E1_r\) and \(E1_{h=0,r}\) were not run. These planned labels
would allow \(r\) to vary in the full class-history candidate and its
association-zero restriction, respectively. Held-out scores intentionally governed this
prespecified continuation decision. They did not select parameters,
admit profile candidates, refine the grid or choose among fits.

\subsection{A descriptive cross-run comparison}

The history-free step-fertility restriction \(E1_{h=0}\) and continuous
\(E0_{\rm free}\) are two different extensions of their matched
fixed-\(r\) \(E0\) baselines. In the same primary windows, both favored
lower relative productivity for large magnitudes during training, and
their forecast differences had the same signs
(Table~\ref{a2:cross-run}). The fixed-\(r\) baseline forecasts agree across
the two implementations to approximately \(5.57\times10^{-5}\) nats
or better over all nine cases, consistent with small local-optimizer
coordinate differences.

% Values reconciled from saved selected training records and forecast CSVs.
\begin{table}[tbp]
\centering\small
\caption{Primary cross-run gains over each study's matched \(E0\)
baseline, in total nats. Training gains are partly ensured by nested
candidate admission. The shared pattern is descriptive and uses the same
historical folds, rather than independent validation.}
\label{a2:cross-run}
\begin{tabular}{rrrrrrr}
\toprule
 & & & \multicolumn{2}{c}{Step fertility} & \multicolumn{2}{c}{Free \(r\)}\\
\cmidrule(lr){4-5}\cmidrule(lr){6-7}
Fold & \(N_{\rm train}\) & \(N_{\rm test}\) & Training & Test & Training & Test\\
\midrule
0 & 301 & 203 & +7.260 & -5.209 & +8.622 & -7.121\\
1 & 504 & 155 & +2.254 & +5.006 & +2.296 & +8.167\\
2 & 659 & 254 & +4.491 & -1.838 & +5.228 & -6.420\\
\bottomrule
\end{tabular}
\end{table}

This correspondence is consistent with both adjustments responding to
weaker apparent training magnitude-productivity scaling. It does not
establish exact equivalence of the fitted features, changing physical
productivity through time, or a specific underprediction of large-event
aftershocks as the cause of the losses. Background, branching and
compensator terms also change. Because the free-\(r\) association/fertility
stage was not executed, the experiments do not determine whether
\(\delta\) would shrink when jointly fitted with \(r\).



\begin{figure}[p]
\centering
\includegraphics[width=\linewidth]{figures/forecast_comparisons.png}
\caption{Saved aggregate forecast differences. Panel A decomposes the
class-history contrast into timing and magnitude terms; diamonds show their
joint difference. Panel B compares free with fixed productivity in each
ground family. Positive values favor the extension. Each setting reuses
the same 612 targets, the panel scales differ, and no uncertainty intervals
are claimed. Source values are the accompanying reconciled gain tables.}
\label{a2:forecast-figure}
\end{figure}
